\section{The mathematical theory implemented by the checker}
\label{sec:unchanged-native-theory}

The reduction in this report concerns the complete validated proof checker.
It does not remove a mathematical axiom, weaken a side condition, or introduce
a new arithmetic principle. For clarity, this section states the existing
theory in full and identifies the inherited proof of its relationship with
iPRA. The source is the ConciseIPRA snapshot at commit
\texttt{42a48d2f04ad72036e8d7e365ea91ea08c76a9e3}.
The attributed foundational proofs are reproduced in the appendices.

The presentation below is the eight-rule calculus called \texttt{IPrf} in
the repository. The existing \texttt{NPrf} kernel also has eight rules, with
weaker equational transport and direct equation-progress introduction;
\path{iprf_iff_nprf} proves that their consequence relations coincide.
The repository additionally has equivalent seven-rule presentations.
None of those mathematical rule counts is improved here.

\subsection{Finite syntax and all primitive rules}

There is one numerical sort. Its terms are
\[
 t ::= x\mid 0\mid S(t)\mid I(n,a;u.h).
\]
Only $u$ is bound in the update body $h$; it is not bound in $n$ or $a$.
Every other parameter of $h$ stays fixed during the iteration. A body is a
fixed finite expression, not a quantified code for a function. In particular,
this syntax supplies no uniform evaluator for arbitrary primitive-recursive
descriptions. All substitutions below avoid capture, and bound variables may
be renamed.

The native formulas are
\[
 A ::= (s=t)\mid Q_yA\mid D_{x,y}A.
\]
The displayed variables are bound in the matrix of their respective binder;
$x$ and $y$ in $D$ are distinct. A context $\Gamma$ is a finite reusable list
of formulas. Write $1:=S0$ and designate $0=1$ as native contradiction.
Neither ordinary implication nor ordinary falsity is primitive syntax here.

Here is the entire primitive inventory. The context in each conclusion is
the same retained context as in its premises.
\begin{align*}
 &\frac{}{\Gamma\vdash A}\quad(A\in\Gamma)
 &&\mathrm{Hyp},\\[4pt]
 &\frac{}{\Gamma\vdash I(0,a;u.h)=a}
 &&I_0,\\[4pt]
 &\frac{}{\Gamma\vdash I(Sn,a;u.h)=h[I(n,a;u.h)/u]}
 &&I_S,\\[4pt]
 &\frac{\Gamma\vdash s=t\qquad\Gamma\vdash F(s)}
        {\Gamma\vdash F(t)}
 &&\mathrm{Tr},\\[4pt]
 &\frac{\Gamma\vdash Q_yB(y)\qquad\Gamma\vdash B(0)}
        {\Gamma\vdash B(t)}
 &&\mathrm{QE},\\[4pt]
 &\frac{\Gamma\vdash D_{x,y}B(x,y)\qquad
         \Gamma,Q_yB(b,y)\vdash A}
        {\Gamma\vdash A}
 &&\mathrm{DE},\\[4pt]
 &\frac{\Gamma,B(t,0)\vdash B(t,Sz)}
        {\Gamma\vdash D_{x,y}B(x,y)}
 &&\mathrm{DI}_{\mathrm{case}},\\[4pt]
 &\frac{\Gamma,E(t,z)\vdash E(t,Sz)}
        {\Gamma\vdash D_{x,y}E(x,y)}
 &&\mathrm{DI}_{\mathrm{progress}}.
\end{align*}
In the last rule $E$ must be a \emph{literal equation}: its outermost
constructor is equality. An equation occurring inside $Q$ or $D$ does not
meet that condition. In Tr, $F(v)$ is any native formula template with one
designated free numerical variable; that variable may occur any finite number
of times, including zero. It is renamed away from the substituted terms and
the retained context. The term $t$ in QE is arbitrary.

In DE, the eigenvariable $b$ is fresh for $\Gamma$, the conclusion $A$,
and the passive parameters of $B$. In either DI rule, $z$ is fresh for
$\Gamma$, the witness $t$, and the passive parameters of the matrix.
The displayed local assumption is discharged. These freshness conditions
prevent a witness from depending on an eigenvariable introduced later.

There are no other primitive equality axioms, successor axioms, induction
schemas, logical connectives, or logical rules. Weakening, contraction,
exchange, fresh renaming, substitution, and assumption cut are derived
operations. For example, if $r=I(0,t;u.0)$, then $I_0$ proves $r=t$;
using that equation as both premises of Tr with $F(v)\equiv(v=t)$ proves
$t=t$. Symmetry, transitivity, and congruence then follow by the usual
equality templates. The full derivations are in
Section~\ref{sec:choice-introduction}.

\subsection{Where the ordinary logical operations come from}

The scalar selector
\[
 C(v,a,b):=I(v,a;u.b),\qquad
 p(y):=I(y,0;u.C(u,1,0))
\]
has $C(0,a,b)=a$ and $C(Sv,a,b)=b$; the update of $C$ ignores $u$.
The term $p$ is the parity iterator, and only its fixed instances
$p(1)=1$ and $p(2)=0$ are needed for the conjunction macro below.

A finite syntactic gate $G_vA$ and selector $M_v(A,B)$ satisfy
\[
 \vdash G_0A,\qquad G_{Sv}A\dashv\vdash A,
 \qquad M_0(A,B)\dashv\vdash A,
 \qquad M_{Sv}(A,B)\dashv\vdash B.
\]
Here interderivability holds under arbitrary additional retained assumptions.
At equations their definitions are
\[
 G_v(s=t):=(C(v,0,s)=C(v,0,t)),
\]
\[
 M_v(s=t,r=q):=(C(v,s,r)=C(v,t,q)).
\]
For a completely finite extension to all formulas, first replace each
$Q_yB$ by the interderivable $D_{r,y}B$, with $r$ unused and fresh.
The gate commutes with $D$. To align selectors of unequal binder depth,
pad an equation $E$ by $D_{r,y}G_yE$; then recurse under the aligned $D$
binders. Section~\ref{subsec:choice-logic} gives the exact terminating
recursion, the padding equivalence, and its proof. This is external recursion
on fixed formula syntax, not an internal truth predicate.

Using those operations, one may take the ordinary logical macros to be
\[
\begin{aligned}
 \bot_*&:=(0=1),\\
 A\mathbin{\to_*}B&:=Q_yM_y(A,B),\\
 A\mathbin{\wedge_*}B&:=Q_yG_yM_{p(y)}(A,B),\\
 \exists_*x\,A(x)&:=D_{x,y}G_yA(x),\\
 A\mathbin{\vee_*}B&:=\exists_*v\,M_v(A,B),\\
 \forall_*x\,A(x)&:=A(0)\mathbin{\wedge_*}Q_xA(x).
\end{aligned}
\]
All new bound variables are fresh; before taking a selector or gate, use
the equivalent $D$-only version just described when necessary.
The equivalent $D$-only inventory in Section~\ref{subsec:choice-logic}
also avoids treating $Q$ as a public formula constructor.

These macros have every intuitionistic natural-deduction rule. Implication
introduction uses the zero selector as its local antecedent and supplies the
successor selector as the consequent; elimination specializes a closure at
$1$. Existential introduction uses its supplied witness, while DE opens a
fresh witness for existential elimination; the zero gate supplies the required
base. Conjunction projections use closure instances $1$ and $2$. Disjunction
uses witnesses $0$ and $1$, followed by the derived numerical case rule for
elimination. A fresh-parameter proof gives universal introduction, and its
zero instance together with its closure gives arbitrary specialization.
Under $0=1$, the selector equations give
$s=C(0,s,t)=C(1,s,t)=t$; external formula induction then proves explosion.
No excluded middle for quantified formulas is introduced.

The resulting numerical case principle and literal-equation induction are
\[
 \frac{\Gamma\vdash A(0)\qquad\Gamma\vdash A(Sz)}
      {\Gamma\vdash A(t)},
 \qquad
 \frac{\Gamma\vdash E(0)\qquad\Gamma,E(z)\vdash E(Sz)}
      {\Gamma\vdash E(t)},
\]
with the usual fresh $z$ and with $E$ a literal equation. In the second
construction, progress DI with an unused witness supplies a closure, and QE
applies it to the given base. The first construction uses case DI and QE.
One must not silently apply the first rule with $z$ free in a retained
assumption. The inherited proofs explicitly discharge that assumption using
the implication macro or use coverage with equality hypotheses instead.

\subsection{The ordinary first-order realization}

To state bi-interpretability in ordinary first-order terms, elaborate each
fixed iterator body $h(u,\vec p)$ to a staged function symbol
$i_h(n,a,\vec p)$ satisfying
\[
 i_h(0,a,\vec p)=a,\qquad
 i_h(Sn,a,\vec p)=h(i_h(n,a,\vec p),\vec p).
\]
Only finitely described bodies are introduced at each stage. This elaboration
changes notation, not the numerical domain.

Let $T_{\mathrm{eq}}$ be intuitionistic first-order logic with actual equality
in this signature, the displayed computation schemas, and
\[
 \forall x\,(x=0\vee\exists y\,x=Sy),
 \qquad (0=S0)\to\bot.
\]
It also has induction for every literal equation $E(n,\vec p)$:
\[
 \bigl(E(0,\vec p)\wedge
       \forall n\,(E(n,\vec p)\to E(Sn,\vec p))\bigr)
       \to\forall n\,E(n,\vec p),
\]
universally closed over the passive parameters. Its contextual induction
rule is the equivalent displayed rule above. The separation sentence is
essential: it connects native contradiction to genuine ordinary falsity.
The equation-only part by itself must not be substituted for this theory.

The ordinary reading $\rho$ fixes terms and equality and is given by
\[
 \rho(Q_yB)=\rho(B(0))\to\forall y\,\rho(B(y)),
\]
\[
 \rho(D_{x,y}B)=
   \exists x\bigl(\rho(B(x,0))\to\forall y\,\rho(B(x,y))\bigr).
\]
Case DI is sound by coverage. Progress DI is sound by equation induction
under a temporarily retained base equation. The other rules have their
ordinary equality or logical meanings. Conversely, the logical macros,
derived numerical cases, and equation induction recover the ordinary
realization. The native translation $\lambda$ of ordinary formulas fixes
equations and replaces the six ordinary logical operations by the macros
above. Sections~\ref{sec:choice-introduction} and~\ref{sec:ordinary-bridge}
prove preservation of sequents and both formula composites, not merely
agreement in the standard natural numbers.

\subsection{The arithmetic bridge to standard iPRA}

Here standard iPRA means intuitionistic first-order arithmetic with one
function symbol for each fixed primitive-recursive description: zero,
successor, projections, composition, and primitive recursion. It has their
defining equations, $Sx\ne0$, successor injectivity, and induction for every
quantifier-free formula, with parameters. The usual intuitionistic logical
rules retain assumptions and enforce fresh variables for universal
introduction and existential elimination.

The inherited arithmetic bridge has a necessary dependency order.
First, a scalar term $P$ is constructed and proved to satisfy
$P0=0$ and $P(Sn)=n$, using no pre-existing pairing or equality detector
(Section~\ref{sec:hybrid-bootstrap}). Next, truncated subtraction and a
Boolean term reflecting actual numerical equality are constructed.
Every fixed quantifier-free formula is then equivalent to its Boolean
characteristic equation, so equation induction yields quantifier-free
induction (Section~\ref{sec:eic-completion}). Only after that step are pairing
and its two projections used to compile general primitive recursion
(Section~\ref{sec:iterator-coding}).

For example, let $\pi,L,H$ satisfy
$L(\pi(a,b))=a$ and $H(\pi(a,b))=b$. A target recursion
\[
 f(0,\vec p)=g(\vec p),\qquad
 f(Sn,\vec p)=h(n,f(n,\vec p),\vec p)
\]
is compiled, after its earlier descriptions $g,h$ have been compiled, by
\[
 \begin{aligned}
 W(u,\vec p)&:=
   \pi\bigl(SL(u),h^*(L(u),H(u),\vec p)\bigr),\\
 V(n,\vec p)&:=I(n,\pi(0,g^*(\vec p));u.W(u,\vec p)),\\
 f^*(n,\vec p)&:=H(V(n,\vec p)).
 \end{aligned}
\]
Equation induction proves $L(V(n,\vec p))=n$, after which the projection
equations prove both target defining equations. The compiler is external
recursion on each finite description. Pairing stores a finite state; the
interpretation domain is not restricted to the image of $\pi$.

In the other direction, iPRA expresses every scalar iterator by ordinary
primitive recursion whose update ignores the recursion counter. It proves
coverage, separation, and equation induction. Thus each theory interprets
the other, and the following stronger statement supplies the required
round trips.

\begin{theorem}[Inherited whole-domain bi-interpretation]
\label{thm:inherited-biinterpretation}
There are effective translations $\Phi$ from native formulas to iPRA and
$\Psi$ from iPRA to native formulas such that, for all finite contexts,
\[
 \Gamma\vdash A
 \quad\Longleftrightarrow\quad
 \Phi\Gamma\vdash_{\mathrm{iPRA}}\Phi A,
\]
\[
 \Delta\vdash_{\mathrm{iPRA}}\varphi
 \quad\Longleftrightarrow\quad
 \Psi\Delta\vdash\Psi\varphi.
\]
Every native $A$ is interderivable with $\Psi\Phi A$, and iPRA proves
$\varphi\leftrightarrow\Phi\Psi\varphi$ for every ordinary formula
$\varphi$. Both interpretations use the entire numerical domain and actual
equality. In each direction the internally definable numerical return
isomorphism has graph $x=y$. In particular,
\[
 \varnothing\vdash 0=S0
 \quad\Longleftrightarrow\quad
 \varnothing\vdash_{\mathrm{iPRA}}\bot.
\]
\end{theorem}

\begin{proof}
Compose the native/ordinary translations $\rho,\lambda$ with the two
arithmetic compilations just described. For a fixed iterator, equation
induction proves uniqueness from its seed and update equation. The same
argument proves uniqueness of a fixed ordinary primitive recursion.
External induction on a finite description therefore shows that its double
translation equals its original value at every internal argument. This is
an internal numerical equality, valid at nonstandard arguments as well as
standard numerals.

The identity numerical graph is consequently total, bijective, and preserves
every interpreted operation. The logical macro equivalences and binder
composites then give both formula round trips by external structural
induction. In particular, $D_{x,y}B$ is interderivable with the native
existential of the native implication from $B(x,0)$ to the native universal
of $B(x,y)$; the analogous closure statement holds for $Q$. The explicit
ordinary separation sentence and native explosion preserve contradiction.
Finite proof translations establish equiconsistency. The detailed
constructions are in Sections~\ref{sec:choice-introduction},
\ref{sec:ordinary-bridge}, and~\ref{sec:iterator-coding}.
\end{proof}

The pinned Lean development records the ordinary arithmetic bridge as
\path{biInterp}, \path{teq_roundTrip}, and
\path{pra_roundTrip}; the native bridge is recorded by
\path{native_iff_pra}, \path{pra_iff_native},
\path{native_roundTrip}, \path{pra_native_roundTrip}, and
\path{native_inconsistent_iff} in
\path{NativeBiInterpretation.lean} and its dependencies.
The comparison \path{iprf_iff_nprf} connects the eight rules stated here
to that kernel. These are attributed existing results; the present
validation refinement neither adds a new Lean theorem nor claims that the
Python implementation has been formally verified in Lean. Its separate
checking argument concerns the representation and validation of certificates
for this unchanged mathematics.
